% ABOUTME: Deep learning setup section of the IEEE report.
% ABOUTME: Sensor windows, padding, normalisation, early stopping and nested engine validation.
\section{Deep Learning Setup}
\label{sec:dl-setup}

\subsection{Sensor windows}
The sequence models read a window of the last $w$ cycles of the 14 useful
sensors and predict the RUL at the window's final cycle. Windows are built
separately for each engine, so no window ever contains cycles from two
engines. We build one window ending at every cycle. This gives the sequence
models exactly the same training and validation rows as the Phase 1 models,
so their cross validation errors are directly comparable. The default length
is $w = 30$: the shortest FD001 test engine has 31 observed cycles, so every
test engine is scored on real history, while the shortest training engine
(128 cycles) is far longer than the window.

\subsection{Padding}
The first $w-1$ cycles of every engine have fewer than $w$ cycles of history.
Rather than drop them, which would remove exactly the healthy early cycles,
we left pad each window with copies of that engine's first cycle. A healthy
engine's sensors are nearly constant, so repeating the first reading adds no
false degradation trend, and the padding never borrows readings from another
engine.

\subsection{Normalisation}
Each sensor is standardised with the mean and standard deviation of the
training engines only. FD001 has a single operating regime, so standardising
by regime and standardising globally coincide. In cross validation the scaler
is refitted inside every fold on that fold's training engines, so no
statistic of a scored engine reaches the model.

\subsection{Training and early stopping}
All sequence models share one training loop: Adam on mean squared error,
mini batches of 256 windows, gradient norms clipped at 1, and seed 42 for
weights and batch order. Training stops after a fixed number of epochs
without improvement in validation loss, and the weights from the best epoch
are restored. Early stopping chooses a hyperparameter, the number of epochs,
by looking at validation error. If it watched the same engines that are
later scored, the reported error would be optimistic. We therefore nest the
split: within each of the five \texttt{GroupKFold} folds, 20\% of the
training engines are held out for early stopping, and the fold's held out
engines are used only for scoring. As in Phase 1, every cycle of the scored
engines is compared with RUL capped at 125.
